\section{Related Work}
\label{sec:related}

We position SSMQ across four research areas and contrast its scheduling guarantees against prior merge-queue and CI optimization designs in Table~\ref{tab:related_taxonomy}.

\begin{table*}[t]
\centering
\footnotesize
\setlength{\tabcolsep}{7pt}
\caption{Qualitative Comparison of Merge Queue \& CI Verification Paradigms.}
\label{tab:related_taxonomy}
\begin{tabular}{lcccc}
\toprule
\textbf{System / Paradigm} & \textbf{Dependency-Aware} & \textbf{Speculation / Batching} & \textbf{Flaky Gating} & \textbf{Fault Attribution} \\
\midrule
Static Batching (Bors / GitHub~\cite{fallahzadeh2023accelerating}) & No (FIFO) & Static batch $K$ & None & $O(\log K)$ binary bisection \\
Linear Trains (GitLab Merge Trains) & No (FIFO) & Linear prefix train & None & Head abort \& restart \\
SubmitQueue (Uber~\cite{ananthanarayanan2019keeping}) & Risk model only & Pass/fail outcome tree & Implicit & Outcome-tree branch split \\
Flake-Aware Bisection~\cite{henderson2023flake} & No (FIFO) & Post-submit culprit search & Bayesian & Noisy binary bisection \\
Predictive Test Selection~\cite{machalica2019predictive,jin2023hybridcisave} & Test granularity & Single-PR build pruning & Orthogonal & N/A (composable) \\
\textbf{SSMQ (Ours)} & \textbf{Cone-Disjoint ($\hat{G}$)} & \textbf{Risk-Stratified Multi-Lane} & \textbf{2-Regime Beta-Binomial} & \textbf{1-Step AST/Stack + Fallback} \\
\bottomrule
\end{tabular}
\end{table*}

\textbf{CI Merge Queues, Batching, and Test Selection.}
Empirical studies of open-source and industrial repositories show that CI adoption improves defect containment and release cadence at the cost of queueing delay and compute contention~\cite{hilton2016usage,zhao2017impact,bernardo2023impact,wessel2023github}. At hyperscale, Google's TAP~\cite{elbaum2014techniques,memon2017taming} introduced batched regression testing, Uber's SubmitQueue~\cite{ananthanarayanan2019keeping} formulated speculative pass/fail outcome trees to sustain monorepo merge throughput, and Meta~\cite{juloori2025scale} decoupled pre-merge and post-merge verification into lean pipelines. Subsequent work investigated failure-history-guided dynamic batch sizing~\cite{bavand2021mining}, parallel batch testing with bisection~\cite{fallahzadeh2023accelerating}, risk-weighted commit bisection~\cite{najafi2019bisecting}, and queue reordering via PR prioritization~\cite{jungwirth2025improving}. A parallel line of work shrinks per-job execution time through test-selection policies based on state transitions~\cite{leong2019assessing}, predictive ML models~\cite{machalica2019predictive}, reinforcement learning~\cite{spieker2017reinforcement}, joint build-and-test filtering~\cite{jin2023hybridcisave}, and pipeline-aware scheduling~\cite{schwendner2025practical}. Whereas test selection prunes test suites \emph{within} a single build and prior batchers group PRs in arrival order, SSMQ intervenes at the \emph{batch formation} layer via warm-up cone disjointness---a mechanism orthogonal to, and directly composable with, per-batch regression test selection.

\textbf{Flaky Tests and Noisy Culprit Isolation.}
Field studies demonstrate that non-deterministic test failures permeate unit, UI, and integration suites~\cite{luo2014empirical,lam2020large,romano2021empirical,tahir2023test} and frequently cluster in correlated infrastructure bursts~\cite{parry2025systemic}. Program-analysis detectors such as DeFlaker~\cite{bell2018laker} and iDFlakies~\cite{lam2019idflakies} pinpoint flaky tests via differential coverage and test-order permutation. On the post-submit side, Ziftci and Reardon~\cite{ziftci2017broke} ranked suspect commits across failing Google builds, while Henderson et al.~\cite{henderson2023flake} formulated Bayesian noisy bisection under flaky outcomes. SSMQ embeds online conjugate Beta-Binomial flakiness estimation directly inside pre-merge queue scheduling and introduces asymmetric thresholds for singletons versus multi-PR batches to balance tail latency against clean-PR false rejections.

\textbf{Semantic Merge Conflicts and Change Impact Analysis.}
Prior work has predicted textual and build conflicts from historical edit patterns~\cite{owhadikareshk2019predicting}, applied PR-level change impact analysis to guide reviewer assignment~\cite{gmen2025enhanced}, and trained neural transformers and language models to synthesize merge resolutions after conflicts surface~\cite{zhang2022trained,svyatkovskiy2022program}. Rather than repairing broken merges post hoc, SSMQ leverages leak-free co-change and directory cones to prevent potentially interfering PRs from co-habiting fast-lane batches while automatically re-verifying isolated PRs when latent interactions occur.

\textbf{Fault Localization and Traceback Attribution.}
Spectrum-based fault localization (SBFL) methods such as Tarantula~\cite{jones2005empirical} and Ochiai~\cite{abreu2009practical} compute statement suspiciousness from per-test coverage matrices, whereas recent approaches fine-tune large language models on mutation-generated stack traces~\cite{jambigi2025fault}. Recording fine-grained coverage matrices on every CI run or invoking multi-second LLM inference inside the scheduler control loop imposes nontrivial operational overhead. SSMQ's \texttt{ASTStackCausalRanker} instead combines zero-instrumentation AST function diffs with 1-hop stack-trace call-graph reachability, isolating suspect PRs in \BenchASTMs{}~ms with fail-safe bisection fallback.
